\section{The Economics of Learning Graph Generation}\label{sec:lg-economics}

A central limitation of any structure-first retrieval architecture is the
up-front cost of constructing the underlying knowledge structure. Section
\ref{sec:disc-limitations} notes that the benchmark's zero-build-cost
assumption for CKG is not strictly accurate: the learning graph itself must
be authored. Historically, this has been the decisive objection to structured
approaches---expert curation is slow, expensive, and does not scale.

This section addresses that objection directly. We first provide a formal
definition of a learning graph, then present a cost model for generating one
via an agentic workflow, and finally extrapolate the trajectory of generation
cost over the next 18--24 months.

\subsection{Formal Definition of a Learning Graph}\label{sec:lg-def}

\begin{definition}[Learning Graph]\label{def:learning-graph}
A \emph{learning graph} is a 4-tuple $G = (C, E, T, \tau)$ where:
\begin{itemize}
  \item $C = \{c_1, c_2, \ldots, c_n\}$ is a finite, non-empty set of
    \emph{concepts}, each a named unit of domain knowledge with a unique
    identifier and a human-readable label.
  \item $E \subseteq C \times C$ is a set of directed \emph{prerequisite
    edges}. An edge $(c_i, c_j) \in E$ asserts that concept $c_i$ must be
    understood before concept $c_j$ can be meaningfully taught or applied.
  \item $T = \{t_1, t_2, \ldots, t_k\}$ is a finite set of \emph{taxonomy
    categories} that partition the concept space into coarse groupings
    (e.g., \texttt{FOUND}, \texttt{CORE}, \texttt{ADV}).
  \item $\tau : C \rightarrow T$ is a total function assigning each concept
    to exactly one taxonomy category.
\end{itemize}
The pair $(C, E)$ must form a \emph{directed acyclic graph} (DAG): there
exists no sequence $c_{i_1}, c_{i_2}, \ldots, c_{i_m}, c_{i_1}$ such that
every consecutive pair is an edge in $E$. Acyclicity ensures a valid
teaching order exists (topological sort over $E$).
\end{definition}

The concrete serialization used throughout this paper and in the McCreary
Intelligent Textbook Corpus encodes $G$ as a four-column CSV:
\texttt{ConceptID}, \texttt{ConceptLabel}, \texttt{Dependencies},
\texttt{TaxonomyID}. The \texttt{Dependencies} column is a pipe-delimited
list of prerequisite \texttt{ConceptID} values, directly representing the
in-edges of each concept in $E$. This serialization is deliberately minimal:
a 200-concept learning graph typically fits in 6--12 KB of plain text.

Three properties follow directly from Definition~\ref{def:learning-graph}
and are load-bearing for the CKG architecture:
\begin{enumerate}
  \item \textbf{Finite, enumerable context.} Because $C$ is finite and the
    CSV serialization is compact, the entire graph fits in an LLM prompt
    for any realistic domain ($|C| \leq 1{,}000$).
  \item \textbf{Deterministic traversal.} Prerequisite chains are computed
    by BFS or DFS over $E$ with no inference required. Query answers for
    structural questions (T2, T3, T4) are functions, not predictions.
  \item \textbf{Closed vocabulary.} The set of valid concept labels is
    exactly $\{\mathrm{label}(c) : c \in C\}$. A retrieval system that
    returns only concepts in $C$ cannot hallucinate entities by construction.
\end{enumerate}

\subsection{Agentic Learning Graph Generation}\label{sec:lg-gen}

The McCreary corpus graphs used in this benchmark were produced with an
agentic workflow (the \texttt{/learning-graph-generator} skill running
under Claude Code) that decomposes generation into three stages:
\begin{enumerate}
  \item \textbf{Concept elicitation.} Given a course description, the
    agent proposes a concept set $C$ of a target size (default $n=200$),
    drawing on the model's domain knowledge and any supplied source
    materials.
  \item \textbf{Dependency assignment.} For each concept $c_j \in C$, the
    agent selects the subset of previously-enumerated concepts on which
    $c_j$ depends, populating $E$ incrementally.
  \item \textbf{Taxonomy assignment and validation.} The agent assigns
    $\tau(c)$ for each $c \in C$, then runs automated validation: cycle
    detection, orphan detection, dependency count distribution, and a
    quality score. Graphs scoring below threshold trigger a correction
    cycle.
\end{enumerate}

A subject-matter expert (SME) reviews the final graph for domain fidelity
and edits edges or labels as needed. In practice, SME review for a
200-concept graph takes 2--4 hours; the agentic generation itself
completes in minutes.

\subsection{Cost Model}\label{sec:lg-cost}

Per-concept generation cost is dominated by two components: output tokens
required to serialize each concept's row ($\sim$200 tokens on average) and
input tokens accumulated as the growing graph must be re-read to avoid
duplicates and cycles ($\sim$500 tokens per concept on average, under a
quadratic-ish accumulation that we approximate as linear for $n \leq 500$).
Fixed overhead covers skill loading, course-description ingestion,
taxonomy-scheme authoring, and the validation pass.

We model generation cost as an affine function of concept count:
\begin{equation}\label{eq:lg-cost}
  \mathrm{Cost}(n) \approx \alpha + \beta \cdot n
\end{equation}
where $\alpha$ is the fixed-overhead cost and $\beta$ is the marginal
per-concept cost. Table~\ref{tab:lg-cost} instantiates
Equation~\ref{eq:lg-cost} for two current-generation Claude models,
calibrated to observed generation sessions on the McCreary corpus.

\begin{table}[htbp]
\centering
\small
\caption{Estimated learning graph generation cost as a function of concept
count, under two current-generation models. Excludes SME review time.
Assumes one validation pass; a correction cycle adds approximately 40\%.}
\label{tab:lg-cost}
\begin{tabular}{lcccc}
\toprule
Model & $\alpha$ (USD) & $\beta$ (USD/concept) & $n=200$ & $n=500$ \\
\midrule
Claude Opus 4.6     & \$1.50 & \$0.025 & \$6.50  & \$14.00 \\
Claude Sonnet 4.6   & \$0.30 & \$0.005 & \$1.30  & \$2.80  \\
\bottomrule
\end{tabular}
\end{table}

Even at Opus-tier pricing, a 200-concept domain graph costs less than ten
dollars of compute to generate. At Sonnet-tier pricing it costs roughly
the price of a cup of coffee. These figures are four to five orders of
magnitude below the cost of traditional expert curation (a team of domain
experts working for weeks or months) that has historically blocked
structure-first retrieval from commercial deployment.

\subsection{Projected Trajectory}\label{sec:lg-projection}

Two compounding trends point toward sharply lower generation cost over
the next 18--24 months:

\begin{enumerate}
  \item \textbf{Declining model pricing at constant capability.} Across
    the preceding two years, intelligence-per-dollar at the Anthropic
    frontier has approximately halved every 9--12 months, as smaller
    models match the capability of their predecessors. Applied to
    Equation~\ref{eq:lg-cost}, both $\alpha$ and $\beta$ scale
    proportionally with per-token pricing.
  \item \textbf{Skill-level efficiency gains.} The current generation
    workflow holds the full emerging graph in context, producing a
    super-linear accumulation term in $\beta$. Decomposing generation
    into cache-efficient passes (concepts first, dependencies batched,
    taxonomy last) collapses the accumulation and has been observed to
    reduce token consumption 2--3$\times$ at any given model tier.
\end{enumerate}

Combining these trends, we project a plausible 2027 cost model of
$\alpha \approx \$0.10$ and $\beta \approx \$0.002$, yielding
approximately \$0.50 for a 200-concept graph---roughly one dollar with
a correction cycle. At that price point, domain graph generation is
effectively free relative to any downstream inference workload.

\subsection{Implications for the Build-Cost Objection}\label{sec:lg-implications}

The traditional objection to structure-first retrieval---\emph{``curating a
domain graph is prohibitively expensive''}---was empirically true for
two decades. It is no longer true. When a 200-concept domain graph can be
generated for dollars of compute and a handful of SME-review hours, the
decision calculus inverts: the question is no longer whether a domain can
afford a learning graph, but whether any high-query-volume domain can
afford \emph{not} to have one, given that per-query retrieval costs fall
by roughly an order of magnitude once the graph exists
(Section~\ref{sec:disc-ckg-wins}).

This shifts where the economic moat lies. Model compute is commoditizing;
the durable value in applying CKG to new domains is the
\emph{SME review loop}---ensuring the generated graph faithfully reflects
the expert consensus of the field---and the \emph{schema design} work
required to extend Definition~\ref{def:learning-graph} beyond
prerequisite-structured domains to domains with richer relation types
(Section~\ref{sec:disc-limitations}).
